\documentclass[11pt,a4paper]{article}
\usepackage[T1]{fontenc}
\usepackage[utf8]{inputenc}
\usepackage[margin=26mm]{geometry}
\usepackage{lmodern,microtype,amsmath,booktabs,tabularx}
\usepackage[hidelinks]{hyperref}
\setlength{\parskip}{0.4em}
\setlength{\parindent}{0pt}
\widowpenalty=10000
\clubpenalty=10000
\title{Reading the BES terms after OCR\\Prospectus closure addendum}
\author{LegalMath development study}
\date{5 October 2026}
\begin{document}
\maketitle

The scanned final terms for three Banco Esp\'irito Santo issues can now be read
by the local prospectus tools. Recovering their text closes an extraction gap,
but it also makes the remaining legal problem more precise: a promise printed
in the final terms must be read with the base prospectus, its amendments and
any applicable external law. The current work supplies reviewed text,
a detailed list of incorporated documents and conditional calculations from
the retained terms. It leaves all three BES issue dossiers unresolved.

ResearchAssistant contains adapters for Marker and MinerU, but the executables
were unavailable in the inspected installation. The existing PyMuPDF installation
provided a working CPU OCR engine after the official English language-data file
was stored locally. No additional executable or Python package was installed.
Seventeen pages from Series 23, 35 and 36 were processed. Eighteen distinct page
images were inspected, including page 287 of the 2010 base prospectus.

\section*{What the page images changed}

OCR inserted an extra character in the Series 23 identifier and damaged the
Series 35 day-count convention. Comparing the text with the retained page images
restored \texttt{PTBEQBOM0010} and \emph{Actual/Actual (ICMA)}. The review records
23 explicit corrections. They also repair issuer names, a date, the ordering of
a clearing-system entry and a damaged qualification concerning Eurosystem
eligibility. Raw OCR, corrected text and the page images remain separately
available, with their identities recorded.

Two empty extractions are genuine blank pages: page 6 of the Series 35 final
terms and page 287 of the 2010 base prospectus. Treating either as a missing
clause would incorrectly block text coverage. Treating every empty extraction
as blank could hide a page of operative terms. The intake
therefore requires an image-bound review for each blank-page declaration.

The review covers the material printed identifiers, dates, payment terms,
exceptions, cross-references and eligibility qualifications. Some immaterial
spacing and signature noise remains. The review neither authenticates signatures
nor certifies a perfect transcription of every glyph. A changed image, missing
review, unsupported text alteration or false OCR declaration causes the new
admission checks to reject the derivative.

All seven retained offering inputs now have complete text coverage under the
recorded blank-page decisions. The reader still returns five abstentions and
two conditional positive readings. The repaired extraction therefore supplies
usable evidence without establishing a new legal answer. The older reader also
reported that OCR had never been performed; the successor admission layer
corrects that field for verified derivatives and retains the original result.

\section*{Which statements prevail}

All fifteen BES bases, supplements and final terms listed in the earlier plan
were already stored locally. The review identifies 29 further incorporated
dependencies, including selected historical conditions and financial statements.
This is a list of programme incorporation requirements, not a conclusion that
every historical condition governs each issue. Series 23 refers to the 2010
base and six supplements. Series 35 refers to the 2013 base and two supplements;
Series 36 adds the February and April 2014 supplements.

The 2013 base incorporates the section on terms and conditions, excluding
undated deeply subordinated notes, from eight earlier base prospectuses. It also
lists precise pages from the financial reports. Reading an entire old prospectus
as operative would extend incorporation beyond the stated section. Conversely,
discarding every old edition because a newer prospectus exists could omit
conditions expressly retained by reference.\footnote{\href{../../prospectus/corner-cases-2026-10-04/sources/jur-more-bes-1774077/source.pdf}{BES base prospectus, 17 July 2013}, printed pp.~58--59, PDF pp.~59--60. The detailed dependency table and exact quotations are retained in \href{DOSSIER-REVIEW.md}{the dossier review}.}

The supplements do not impose a universal \emph{latest document wins} rule.
Their precedence language concerns inconsistent statements or modification to
the extent applicable. Some amendments are explicit. The February 2014
supplement replaces the Portuguese tax section and deletes the former tax
forms. The April supplement replaces Summary B12 and the statements about
significant or material change. Its incorporated 2013 accounts were still
subject to the scheduled shareholder approval when the supplement was
published; the later occurrence of that approval needs separate evidence.

Section numbering reveals a concrete ambiguity. The November 2013 supplement
renumbers the financial-transaction-tax risk as 1.17. The April 2014 supplement
inserts a KPMG audit risk under the same number. The successor retains both
records and reports the conflict. It does not let the second numerical key erase
the first risk.\footnote{\href{../../prospectus/corner-cases-2026-10-04/sources/jur-more-bes-1843910/source.pdf}{Supplement, 6 November 2013}, p.~2; \href{../../prospectus/corner-cases-2026-10-04/sources/jur-more-bes-2001058/source.pdf}{Supplement, 23 April 2014}, p.~6.}

\section*{Calculating a conditional coupon}

Series 35 specifies an annual rate of 4 percent and an issue amount of
EUR~750 million. Series 36 specifies 2.625 percent on the same issue amount.
Both select Actual/Actual (ICMA). For Interbolsa notes, the base prospectus
applies the interest rate to the full outstanding nominal amount. Let \(N\)
denote that amount in euros, \(r\) the annual rate as a decimal, \(d\) the
accrual days, \(D\) the days in the relevant determination period and \(m\)
the number of determination dates per year. For the short or regular period
described in Condition 5(a), the unrounded coupon is
\[
 I=Nr\,\frac{d}{Dm}.
\]
Dates include the start and exclude the end. For a long period spanning the
two determination periods covered by the printed formula, the fraction becomes
\(d_1/(D_1m)+d_2/(D_2m)\).\footnote{\href{../../prospectus/corner-cases-2026-10-04/sources/jur-more-bes-1774077/source.pdf}{BES base prospectus}, Condition 5(a), printed pp.~95--96, PDF pp.~96--97. The implementation retains exact rational arithmetic until currency rounding.}

For a full annual period, \(d=D\) and \(m=1\). If the original nominal amount
remains outstanding, the coupons are therefore EUR~30 million for Series 35
and EUR~19,687,500 for Series 36. These amounts are conditional calculations,
not evidence of what holders were actually paid. Early redemption, purchases,
cancellation, external law and the complete issue documentation still matter.

A partial period exposes why the denominator cannot silently become 365.
From 21 January to 21 July 2016 there are 182 accrual days in a 366-day annual
determination period. On EUR~100,000 at 4 percent, the printed formula gives
\[
 100{,}000(0.04)\frac{182}{366}=1{,}989.071\ldots,
\]
or EUR~1,989.07 under the printed half-up cent rule. The base also allows an
applicable market convention. The calculator requires the caller to select
the supported convention explicitly; an unknown convention produces no amount.
For Interbolsa notes, rounding occurs after calculating on the outstanding
nominal amount. Rounding each holding before aggregation can give a different
answer at half-cent boundaries.

\section*{Qualifications that must survive}

The final terms deem a clearing-system notice delivered on the second business
day after it is given to Interbolsa. The calculator needs a complete calendar
for those dates. It cannot infer holidays from weekdays. Likewise, checking
the tax-redemption notice's 30--60-day interval does not establish the tax
trigger or settle a disputed legal counting rule.

The final terms answer \emph{Yes} to intended arrangements for Eurosystem
eligibility, then explain that this does not necessarily mean recognition as
eligible collateral. Recognition depends on the eligibility criteria.
The successor records the intended arrangement and leaves actual collateral
eligibility unknown.\footnote{Series 35 \href{../../prospectus/corner-cases-2026-10-04/sources/jur-more-bes-1922255/source.pdf}{final terms}, pp.~2 and 5; Series 36 \href{../../prospectus/jurisdiction-2026-10-04/sources/jur-bes-2014-ptbeqkom0019/source.pdf}{final terms}, pp.~2 and 5.}

The focused suite passes 66 tests covering derivative provenance, altered
negation, blank pages, day counts, rounding, missing facts and section conflicts.
The preserved 4 October candidate has its own 885-test result. Neither count
measures independent legal accuracy. Twenty-five review forms are prepared
for qualified adjudicators, and the original cohort's 1,467 unresolved records
remain available across 1,307 source spans. Those records are a development
inventory, not an error-rate estimate.

Further closure requires exact incorporated reports and operative instruments,
resolution of document conflicts, source-supported repairs across the original
cohort, actual investigation facts and independent labels. A future unseen
evaluation must exclude every prospectus already inspected. Human review of
this rendered explanation and intended-use acceptance also remain pending.
The present repairs can be retained and tested without treating those pending
steps as completed.
\end{document}
